\section{Discussion}
\todo[inline]{Interpret the results. Relate findings to your original objectives. Discuss limitations and possible improvements.}

We observe best tuned model, XGBoost has moderate performance on the test data, with having slightly better
precision (0.658) than recall (0.595). It means for this test dataset when the model has more false negatives 
than false positives. Since the ultimate aim of the model application is to prioritize new drug-disease indications,
better precision yields makes results that are predicted as positive more reliable, but 
worse recall means that some combination may be (potentially systematically and unfairly) ignored.
Other models performed slightly worse, especially kNN, this suggests that those
models didn't capture well non-linear patterns in noisy multidimensional data or failed to 
handle class imbalance. 

Heterogenous GNN was trained and tested separately and yielded a large number
of false positives resulting in low precision of 0.429. This poor performance can be explained by
poor handling of class imbalance on our side, since we used binary cross entropy loss with positive weights 
equal to number of times success labels are less prevalent and it might have been to stringent to the model.
We should further analyze the correct ways to handle class imbalance in GNN including proper hyperparameter and classification
threshold tuning. But considering other potential explanations for a poor performance, 
we could argue that the size of training data was not enough to properly train such a complex
model. We also need to discover ways to handle potential overfitting issues with GNN.


Two most important features identified by permutation method are pathway similarity and number of indications for a disease,
and both these variables need to be high for assignment of large success probability. It is expected to see pathway similarity
being predictive since this metric captures shared molecular mechanisms behind action of a drug and pathogenesis of a disease.
It is also unfortunately expected to see number of indications being predictive of success since this feature reflects
how intensively researched is the disease, this constitutes a bias towards well-studied diseases and potentially neglecting 
rare or poorly studied ones. However, analysis of feature distributions for correctly classified and misclassified cases didn't reveal any prominent bias.
So probably the importance of number of indications reflects a bias in clinical research itself rather than a bias in the model.

We can also identify


... biases: publication, data availability, rare diseases, a lot of unknowns, obvious negatives are not tried ...



